%=====================================================================
%  Seminar Report — Verification Methodologies for RISC-V Processors
%  Kevin Jose (U2303131)  |  Dept. of CSE, RSET  |  September 2026
%  Format mirrors the CSE 2026 template and the Kavitha reference report.
%---------------------------------------------------------------------
%  Compile: pdflatex seminar_report.tex  (run twice for ToC/refs)
%  Figures: replace \figplaceholder{...} with
%           \includegraphics[width=0.8\textwidth]{filename}
%=====================================================================
\documentclass[12pt,a4paper]{report}

\usepackage[utf8]{inputenc}
\usepackage[T1]{fontenc}
% Default Computer Modern serif font (matches the reference report)
\usepackage[a4paper,left=1.25in,right=1in,top=1in,bottom=1in]{geometry}
\usepackage{setspace}
\usepackage{graphicx}
\usepackage{amsmath}
\usepackage{amssymb}      % \mathfrak (title heading), \checkmark
\usepackage{booktabs}
\usepackage{array}
\usepackage{longtable}
\usepackage{tabularx}
\usepackage{enumitem}
\usepackage[hidelinks]{hyperref}
\usepackage{fancyhdr}
\usepackage{caption}
\usepackage{pifont}       % \ding for tick/cross marks in the matrix

\newcommand{\yes}{\ding{51}}   % ✓
\newcommand{\no}{\ding{55}}    % ✗

\onehalfspacing
\captionsetup{font=small,labelfont=bf,labelsep=colon}
\renewcommand{\arraystretch}{1.3}

% caption format "Figure X.Y: ..." and "Table X.Y: ..."
\captionsetup[figure]{name=Figure}
\captionsetup[table]{name=Table}

% ---- placeholder box for figures to be inserted later ----
\newcommand{\figplaceholder}[1]{%
  \fbox{\parbox[c][4.2cm][c]{0.78\textwidth}{\centering \itshape
  [ Insert figure here: #1 ]}}}

\pagestyle{fancy}
\fancyhf{}
\fancyhead[R]{\thepage}
\renewcommand{\headrulewidth}{0pt}

\begin{document}

%=====================================================================
% TITLE PAGE
%=====================================================================
\begin{titlepage}
\centering
\vspace*{0.3cm}
\IfFileExists{rset_logo.png}{\includegraphics[width=2.4cm]{rset_logo.png}\\[0.5cm]}{}
{\large\itshape Seminar Report On\par}
\vspace{0.9cm}
{\LARGE\bfseries Verification Methodologies for RISC-V Processors\par}
\vspace{1.0cm}
{\itshape Submitted in partial fulfillment of the requirements for the\\
award of the degree of\par}
\vspace{0.7cm}
{\LARGE $\mathfrak{Bachelor}\;\mathfrak{of}\;\mathfrak{Technology}$\par}
\vspace{0.3cm}
{\itshape in\par}
\vspace{0.3cm}
{\large\bfseries Computer Science and Engineering\par}
\vspace{0.9cm}
{By\par}
\vspace{0.4cm}
{\large Kevin Jose (U2303131)\par}
\vspace{0.9cm}
{Under the guidance of\par}
\vspace{0.4cm}
{\large Dr. Jomina John\par}
\vspace{1.1cm}
{\bfseries Department of Computer Science and Engineering\\
Rajagiri School of Engineering \& Technology (Autonomous)\par}
{(Parent University: APJ Abdul Kalam Technological University)\\
Rajagiri Valley, Kakkanad, Kochi, 682039\\
September 2026\par}
\end{titlepage}

%=====================================================================
% CERTIFICATE
%=====================================================================
\thispagestyle{empty}
\begin{center}
{\large\bfseries CERTIFICATE}
\end{center}
\vspace{1.0cm}

This is to certify that the seminar report entitled ''\textbf{Verification
Methodologies for RISC-V Processors}'' is a bonafide record of the work done by
\textbf{Kevin Jose (U2303131)}, submitted to Rajagiri School of Engineering \&
Technology (RSET) (Autonomous) in partial fulfillment of the requirements for the
award of the degree of Bachelor of Technology (B.\,Tech.) in Computer Science and
Engineering during the academic year 2026--2027.

\vspace{2.4cm}
\noindent
\begin{minipage}[t]{0.5\textwidth}
Dr. Jomina John\\
Associate Professor\\
Dept. of CSE\\
RSET
\end{minipage}%
\begin{minipage}[t]{0.5\textwidth}
\raggedleft
Dr. Jincy J. Fernandez\\
Associate Professor\\
Dept. of CSE\\
RSET
\end{minipage}

\vspace{2.2cm}
\begin{center}
Dr. Mary Priya Sebastian\\
Associate Professor \& HoD\\
Dept. of CSE\\
RSET
\end{center}

%=====================================================================
% FRONT MATTER
%=====================================================================
\clearpage
\pagenumbering{roman}
\setcounter{page}{1}

% ---- ACKNOWLEDGMENT ----
\chapter*{ACKNOWLEDGMENT}
\addcontentsline{toc}{chapter}{Acknowledgment}
\thispagestyle{fancy}

\hspace{0.5cm} I wish to express my sincere gratitude towards Rev.\ Dr.\ Jaison Paul
Mulerikkal CMI, Principal of RSET, and Dr.\ Mary Priya Sebastian, Associate Professor
\& Head of the Department of Computer Science and Engineering for providing me with the
opportunity to undertake my seminar, ''\textit{Verification Methodologies for RISC-V
Processors}''.

\vspace{0.4cm}
\hspace{0.5cm} I am highly indebted to my seminar coordinators, Dr.\ Jincy J. Fernandez,
Associate Professor and Mr.\ Ajith S., Assistant Professor, Department of CSE, for their
valuable support.

\vspace{0.4cm}
\hspace{0.5cm} It is indeed my pleasure and a moment of satisfaction for me to express my
sincere gratitude to my seminar guide, Dr.\ Jomina John, Associate Professor, Department
of CSE, for her patience and all the priceless advice and wisdom she has shared with me.

\vspace{0.4cm}
\hspace{0.5cm} Last but not the least, I would like to express my sincere gratitude
towards all other teachers and friends for their continuous support and constructive
ideas.

\vspace{1.0cm}
\begin{flushright}
Kevin Jose
\end{flushright}

% ---- ABSTRACT ----
\chapter*{Abstract}
\addcontentsline{toc}{chapter}{Abstract}

\hspace{0.5cm} RISC-V is an open, royalty-free instruction set architecture (ISA) that
has been widely adopted across academia and industry, largely because it permits
designers to add their own custom instructions. However, once a processor is fabricated a
functional bug can no longer be patched, which makes pre-silicon verification critical.
The official RISC-V compliance suite is handwritten and implementation-independent: it
applies the same tests to every core and therefore cannot systematically target a specific
pipeline depth, hazard logic, or cache configuration. Consequently, passing the compliance
suite is not equivalent to being correct.

This seminar presents a study of verification methodologies for RISC-V processors through
a review of four recent journal publications and an in-depth analysis of the selected base
paper. The four supporting papers span the principal verification paradigms: high-level
formal verification through model checking ($\chi$RVFormal), golden-model functional
verification of a vector accelerator using a UVM environment, coverage-guided RTL fuzzing
(INSTILLER), and polynomial formal verification using Binary Decision Diagrams (PFV). The
base paper, PATARA, proposes a self-testing framework built on the REVERSI principle, in
which an operation and its inverse are applied so that a test checks itself without any
golden reference model, together with a twin-based method for verifying non-invertible
operations against an already-verified core.

The report discusses the methodology, system architecture, experimental setup, results,
and comparative analysis of the base paper, and identifies the recurring theme across the
studied literature: the two independent axes of verification, namely where test stimuli
come from and how correctness is decided. The study is directly relevant to the author's
project, which implements custom AI-acceleration instructions on the Shakti C-Class core,
and it establishes a practical verification path for custom instructions that, by
definition, have no golden reference.

% ---- Lists ----
\clearpage
\addcontentsline{toc}{chapter}{List of Abbreviations}
\chapter*{List of Abbreviations}
\begin{longtable}{ll}
ALU   & Arithmetic Logic Unit\\
BDD   & Binary Decision Diagram\\
BMC   & Bounded Model Checking\\
BSV   & Bluespec SystemVerilog\\
CI    & Continuous Integration\\
CSR   & Control and Status Register\\
DUT   & Design Under Test\\
FSM   & Finite State Machine\\
HDL   & Hardware Description Language\\
ISA   & Instruction Set Architecture\\
ISS   & Instruction Set Simulator\\
OVI   & Open Vector Interface\\
RTL   & Register Transfer Level\\
RVV   & RISC-V Vector Extension\\
SIMD  & Single Instruction, Multiple Data\\
SMT   & Satisfiability Modulo Theories\\
SVA   & SystemVerilog Assertion\\
UVM   & Universal Verification Methodology\\
VACO  & Vector Ant Colony Optimization\\
VPU   & Vector Processing Unit\\
\end{longtable}

\clearpage
\addcontentsline{toc}{chapter}{List of Figures}
\listoffigures

\clearpage
\addcontentsline{toc}{chapter}{List of Tables}
\listoftables

% ---- CONTENTS ----
\clearpage
\tableofcontents

%=====================================================================
% MAIN MATTER
%=====================================================================
\clearpage
\pagenumbering{arabic}
\setcounter{page}{1}

%---------------------------------------------------------------------
\chapter{Introduction}

\section{Background}
Modern computing increasingly depends on the ability to design processors tailored to a
specific workload. The RISC-V instruction set architecture has accelerated this trend by
providing an open, royalty-free, and modular ISA that anyone can implement or extend. A
distinctive feature of RISC-V is its reserved encoding space for custom instructions, which
allows designers to add domain-specific operations---for example, instructions that
accelerate machine-learning or signal-processing workloads at the edge. This flexibility,
however, shifts a significant burden onto verification: a custom processor is only useful if
it can be shown to behave correctly, and once a chip is fabricated a functional error can
no longer be corrected~\cite{riscvman}.

\section{Overview of the Seminar Topic}
Verification is the process of establishing that a processor implementation behaves exactly
as its specification requires. For RISC-V cores this is complicated by two factors. First,
the official compliance suite is handwritten and implementation-independent, so it cannot
target the pipeline depth, forwarding logic, or cache organisation of any particular core.
Second, most automated verification techniques rely on a \emph{golden reference model}---a
trusted implementation such as the Spike simulator---to decide whether an output is
correct; but a custom instruction that a designer has invented is, by definition, absent
from any such reference. This seminar surveys the principal methodologies that address
these challenges: formal verification, golden-model functional verification, coverage-guided
fuzzing, and self-testing.

\section{Objectives of the Study}
The main objectives of this seminar are:
\begin{itemize}[nosep]
  \item To understand the major methodologies used to verify RISC-V processor
        implementations.
  \item To analyse the working principle, strengths, and limitations of each
        verification approach.
  \item To perform a detailed study of the selected base paper on self-testing
        verification.
  \item To examine how correctness is established for a custom RISC-V core and
        co-processor without a golden reference model.
  \item To identify how these methodologies apply to the verification of custom
        instructions in the proposed project.
\end{itemize}

\section{Scope of the Study}
The study covers five peer-reviewed journal papers published between 2023 and 2026, each
representing a distinct verification approach. It discusses the underlying principles,
experimental results, advantages, and limitations of each method, with particular emphasis
on how correctness is decided and where test stimuli originate. Detailed
electronic-design-automation tool internals, complete formal-logic proofs, and the full
micro-architecture of every evaluated core are beyond the scope of this report; the focus
is on the verification methodology and its applicability to custom instruction extensions.

\section{Organization of the Report}
Chapter~1 introduces the domain, objectives, and scope. Chapter~2 reviews the four
supporting papers, compares them, identifies the research gaps, and motivates the base
paper. Chapter~3 analyses the base paper in detail---its methodology, architecture,
experimental setup, results, and limitations. Chapter~4 summarises the key findings,
relates them to the proposed project, and presents the overall conclusion.

%---------------------------------------------------------------------
\chapter{Literature Review}

\section{Overview of Existing Research}
Processor verification has undergone significant development as designs have grown more
complex and as open ISAs such as RISC-V have made custom cores widespread. Broadly, the
field can be understood along two independent axes: \emph{where test stimuli come from}
(systematic enumeration, constrained-random generation, coverage-guided feedback, or none
at all in the case of exhaustive formal methods), and \emph{how correctness is decided}
(comparison against a golden model, a self-checking property, or a mathematical proof).

Early verification relied largely on directed, hand-written tests and on the official
compliance suite. While simple to apply, such tests are implementation-independent and
cannot target the timing-dependent corner cases of a specific pipeline. Subsequent work
introduced constrained-random stimulus generators checked against a golden reference model
such as the Spike instruction-set simulator, and, more recently, coverage-guided fuzzing
that uses feedback to steer test generation. In parallel, formal verification matured from
verifying small blocks to proving properties of complete processors, using engines such as
model checking and Binary Decision Diagrams.

The four papers reviewed in this chapter occupy different positions on the two axes and
therefore frame the position of the base paper. $\chi$RVFormal uses formal model checking
against a compact reference model; the RISC-V Vector Accelerator paper uses
constrained-random stimuli checked against the Spike golden model; INSTILLER uses
coverage-guided fuzzing; and PFV uses BDD-based formal verification with a proven bound on
verification time. Together they motivate the selection of the base paper, which uniquely
avoids a golden model altogether.

%........................ Paper 1 ....................................
\section{Review of Research Paper 1}
\noindent\textbf{Paper Title:} $\chi$RVFormal: Formal Verification of RISC-V Processor
Chisel Designs\\
\textbf{Authors:} S. Shen, S. Chen, Y. Liu, L. Zhang, F. Song, and Z. Wu\\
\textbf{Year:} 2026

\subsection{Objective}
The objective of this research is to make formal verification of RISC-V processors
practical by replacing the large, hard-to-maintain reference model used by existing tools
with a compact, high-level model that additionally covers privileged instructions and
virtual memory, and by enabling the verification of pipelined and out-of-order cores.

\subsection{Methodology}
The authors develop a modular, parameterised RISC-V ISA reference model written in Chisel,
a high-level hardware description language, in approximately 4{,}609 lines, compared with
the roughly 104{,}911 lines of Verilog used by the existing riscv-formal tool. Because the
reference model executes one instruction per step while a real processor is pipelined and
may execute out of order, a synchronisation mechanism aligns the two at each instruction
commit. Two strategies are provided: StateCheck, which compares the full architectural
state at each commit and suits in-order cores, and ActionCheck, which compares each
instruction's individual effect and scales to out-of-order cores. The resulting transition
system is discharged by automated model checkers.

\subsection{Key Contributions}
\begin{itemize}[nosep]
  \item A compact Chisel reference model covering unprivileged and privileged
        instructions and Sv39 virtual memory.
  \item A synchronisation mechanism (StateCheck and ActionCheck) enabling comparison
        against pipelined and out-of-order designs.
  \item Discovery of previously unknown, designer-confirmed bugs in open-source cores.
\end{itemize}

\subsection{Advantages}
\begin{itemize}[nosep]
  \item Proves conformance to the specification for all inputs rather than for
        sampled tests.
  \item The reference model is small and maintainable.
  \item Scales to a complex out-of-order core (BOOM).
\end{itemize}

\subsection{Limitations}
\begin{itemize}[nosep]
  \item The reference model and synchronisation mechanism are not themselves formally
        verified.
  \item Out-of-order execution combined with virtual-memory operations is not fully
        supported.
  \item Evaluation is limited to three open-source processors, and the reference model
        covers only a practical subset of the ISA.
\end{itemize}

%........................ Paper 2 ....................................
\section{Review of Research Paper 2}
\noindent\textbf{Paper Title:} Functional Verification of a RISC-V Vector Accelerator\\
\textbf{Authors:} V. Jim\'enez, M. Rodr\'iguez, M. Dom\'inguez, et al.\\
\textbf{Year:} 2023

\subsection{Objective}
The objective is to functionally verify Vitruvius+, an industrial RISC-V vector
accelerator connected to a scalar core over the Open Vector Interface, to a level of
confidence sufficient for silicon tape-out.

\subsection{Methodology}
A reusable Universal Verification Methodology (UVM) environment is built around the
seven-channel Open Vector Interface, with one agent per sub-interface; only the ISSUE
channel is randomised while the others react. Constrained-random vector instructions are
generated using an extended version of Google's RISCV-DV generator, and the Spike
simulator serves as the golden reference model, with results compared through a UVM
scoreboard. More than fifty SystemVerilog assertions, concentrated on the memory
interfaces, provide observability, and continuous-integration regressions run for
approximately one year.

\subsection{Key Contributions}
\begin{itemize}[nosep]
  \item An industrial demonstration of golden-model functional verification of a
        vector accelerator.
  \item An extension of the RISCV-DV generator to produce vector instructions.
  \item A detailed account of the coverage achieved and the classes of errors found.
\end{itemize}

\subsection{Advantages}
\begin{itemize}[nosep]
  \item Realistic, industry-standard methodology.
  \item Validated by silicon tape-out.
  \item High functional coverage (95.79\%).
  \item Strong observability from interface assertions.
\end{itemize}

\subsection{Limitations}
\begin{itemize}[nosep]
  \item Verification is mainly at the interface level rather than of individual VPU
        submodules.
  \item Code coverage, and especially toggle coverage, remained low.
  \item The multi-agent UVM environment was difficult to maintain.
  \item The work targets the deprecated RVV~0.7.1 draft rather than the ratified
        specification.
\end{itemize}

%........................ Paper 3 ....................................
\section{Review of Research Paper 3}
\noindent\textbf{Paper Title:} INSTILLER: Toward Efficient and Realistic RTL Fuzzing\\
\textbf{Authors:} G. Zhang, P. Wang, T. Yue, D. Liu, Y. Guo, and K. Lu\\
\textbf{Year:} 2024

\subsection{Objective}
The objective is to make coverage-guided RTL fuzzing of processors more efficient and more
realistic by keeping test inputs short and by exercising interrupts and exceptions that
prior fuzzers ignored.

\subsection{Methodology}
Random instruction streams are executed on both an instruction-set simulator and the RTL,
and any disagreement (a ``mismatch'') is flagged as a candidate bug; coverage feedback
steers subsequent generation. An input-distillation technique based on ant colony
optimisation (VACO) keeps test sequences short while preserving coverage: coverage points
are treated as ``cities'' and candidate test sequences as ``ants'', with pheromone
reinforcing the instructions that reach coverage efficiently. Realistic multi-interrupt and
exception injection with priorities, and hardware-aware seed selection and mutation, further
improve effectiveness.

\subsection{Key Contributions}
\begin{itemize}[nosep]
  \item The VACO input-distillation algorithm based on ant colony optimisation.
  \item Realistic multi-interrupt and exception handling with priorities.
  \item Hardware-aware seed selection and mutation.
\end{itemize}

\subsection{Advantages}
\begin{itemize}[nosep]
  \item Achieves 29.4\% higher coverage and 17.0\% more mismatches than DiFuzzRTL.
  \item Makes test inputs 79.3\% shorter, improving simulation efficiency.
  \item An ablation confirms that removing VACO lengthens inputs by 417\%.
\end{itemize}

\subsection{Limitations}
\begin{itemize}[nosep]
  \item INSTILLER does not outperform DiFuzzRTL on every target.
  \item The evaluation compares mainly against a single baseline and uses four CPU
        cores, two of which are OpenRISC rather than RISC-V.
  \item VACO with relationship extraction adds overhead that offsets part of the
        benefit of shorter inputs.
\end{itemize}

%........................ Paper 4 ....................................
\section{Review of Research Paper 4}
\noindent\textbf{Paper Title:} Polynomial Formal Verification of a RISC-V Processor\\
\textbf{Authors:} L. Weingarten, K. Datta, and R. Drechsler\\
\textbf{Year:} 2025

\subsection{Objective}
The objective is to make formal verification predictable by proving not only that a
processor is correct but also a mathematical bound on the time the verification itself
requires.

\subsection{Methodology}
The design is verified using Binary Decision Diagrams, a canonical representation of logic
in which two equivalent circuits produce identical graphs, so equivalence checking reduces
to a pointer comparison. Each instruction is isolated by pinning the control signals that
select it, marking the data inputs as symbolic, and deleting nodes that collapse to
constants; the remaining logic is compared against a reference BDD built from the
instruction's mathematical definition. The control unit is verified as a finite-state-machine
walk. The authors then prove that the BDD sizes---and hence the verification time---grow
polynomially with the datapath width.

\subsection{Key Contributions}
\begin{itemize}[nosep]
  \item Proven polynomial complexity bounds for verification time (e.g.\ $O(n)$ for
        fetch and $O(n^2)$ for execute/decode).
  \item A complete correctness guarantee over all inputs.
  \item Verification of all 32 control signals, up from 15 in prior work.
\end{itemize}

\subsection{Advantages}
\begin{itemize}[nosep]
  \item Provides a complete correctness guarantee.
  \item Offers a predictable, polynomial bound on verification time.
  \item Addresses the unpredictable-runtime concern that limits industrial adoption of
        formal methods.
\end{itemize}

\subsection{Limitations}
\begin{itemize}[nosep]
  \item The evaluated core (MicroRV32) is multi-cycle and non-pipelined, with no
        hazards or forwarding.
  \item Interrupts are excluded from the scope.
  \item The BDD approach does not scale to multiply, divide, or floating-point units,
        which grow exponentially.
\end{itemize}

%........................ Comparative analysis ......................
\section{Comparative Analysis}
The four reviewed papers, together with the base paper, demonstrate the range of
verification methodologies available for RISC-V processors. $\chi$RVFormal and PFV apply
formal methods and therefore reason about all inputs, but the former bounds its proof to a
fixed cycle depth while the latter applies only to a trivial, non-pipelined core. The
Vector Accelerator paper and INSTILLER scale to realistic designs, but both check only
sampled inputs and depend on a reference---Spike or an instruction-set simulator---that
describes the standard ISA. Although each paper makes valuable contributions, none of the
four supporting papers can verify a custom instruction without its reference first being
extended by hand.

Table~\ref{tab:compare} compares the five papers across the properties most relevant to
the verification of a custom RISC-V core. As with the base paper of any well-motivated
survey, the selected base paper (Paper~V) is the only one that satisfies the
distinguishing requirements: it needs no golden reference model, it verifies custom
instructions directly, and it remains usable post-silicon.

\begin{table}[h]
\centering
\caption{Comparative Analysis of the Selected Papers}
\label{tab:compare}
\small
\begin{tabularx}{\textwidth}{Xccccc}
\toprule
\textbf{Property} & \textbf{Paper I} & \textbf{Paper II} & \textbf{Paper III}
& \textbf{Paper IV} & \textbf{Paper V (Base)}\\
\midrule
No golden reference model      & \no  & \no  & \no  & \no  & \yes\\
Formal proof over all inputs   & \yes & \no  & \no  & \yes & \no\\
Detects real design bugs       & \yes & \yes & \yes & \no  & \yes\\
Scales to pipelined / OoO core & \yes & \yes & \yes & \no  & \yes\\
Verifies custom instructions   & \no  & \no  & \no  & \no  & \yes\\
Usable post-silicon            & \no  & \no  & \no  & \no  & \yes\\
\bottomrule
\end{tabularx}\\[4pt]
\footnotesize Paper~I: $\chi$RVFormal; Paper~II: Vector Accelerator; Paper~III: INSTILLER;
Paper~IV: PFV; Paper~V: PATARA (base paper).
\end{table}

\section{Research Gaps Identified}
The literature review reveals the following gaps in existing RISC-V verification methods:
\begin{itemize}[nosep]
  \item Most methods depend on a reference---a formal model, the Spike simulator, or an
        instruction-set simulator---that describes only the standard ISA.
  \item None of the supporting methods can verify a custom, designer-defined instruction
        without extending its reference by hand.
  \item Methods that give the strongest guarantees (formal verification) either bound
        their proofs to a fixed depth or apply only to trivial, non-pipelined cores.
  \item Methods that scale to real designs (golden-model testing and fuzzing) check only
        sampled inputs, confirming correctness only for the tests actually run.
  \item Post-silicon validation is difficult for methods that require an external
        reference model.
\end{itemize}

\section{Motivation of the Base Paper}
The base paper, PATARA, was selected because it directly addresses the most important of
these gaps. It verifies a custom RISC-V core and its co-processor \emph{without} any golden
reference model: invertible instructions check themselves through the REVERSI principle,
and non-invertible operations are checked by a twin method against an already-verified core.
Because it targets custom instructions that have no reference---precisely the situation of
the author's project---and because it demonstrates concretely that passing the compliance
suite is not the same as being correct, it forms an ideal basis for detailed study.

\vspace{0.3cm}
\noindent\textbf{Chapter Conclusion}\\[2pt]
This chapter presented a detailed review of four verification methodologies for RISC-V
processors---formal model checking, golden-model functional verification, coverage-guided
fuzzing, and polynomial formal verification. The comparative analysis positioned the papers
along the two verification axes and identified a common limitation: every reviewed method
depends on a reference for the standard ISA and therefore cannot verify a custom instruction
directly. Based on these gaps, the base paper was found to be highly relevant because it
establishes correctness without any golden reference model. The findings of this review
provide the motivation for the detailed analysis of the base paper presented in the next
chapter.

%---------------------------------------------------------------------
\chapter{Detailed Analysis of the Base Paper}

\section{Overview of the Base Paper}
The base paper selected for this seminar is ''A Self-Testing Framework for Verification and
Validation of a RISC-V-Based System with a Co-processor'' by Gesper et al., published in the
International Journal of Parallel Programming (2026). The paper presents PATARA, a framework
that automatically generates test programs that check themselves, so that a custom RISC-V
core and an attached vector co-processor can be verified without any external golden
reference model.

The central finding of the paper is that the official RISC-V compliance suite passed the
authors' processor with zero detected bugs while leaving roughly a fifth of the design's
logical conditions untested, whereas better, implementation-aware tests reached full
coverage and revealed real bugs in the same ``passing'' design. This directly challenges the
common assumption that passing the compliance suite is equivalent to correctness. A major
contribution of the work is the twin-based method, which extends self-testing to operations
that cannot be inverted by comparing them against an already-verified core.

\section{Proposed Methodology}
The methodology proposed in the base paper combines a self-checking principle, an automated
generation tool, a set of extensions targeting pipeline corner cases, and a twin method for
operations that cannot be inverted. The complete flow converts an ISA description into
self-checking assembly tests that require no external reference.

\subsection{The REVERSI Self-Testing Principle}
Conventional verification asks whether an output equals the correct value, which forces the
correct value to be computed elsewhere by a golden reference model. REVERSI avoids this: an
operation is applied and then its inverse is applied, and the result must return to the
original value. Adding a number and then subtracting it must restore the starting value; if
it does not, the hardware is faulty---and the correct intermediate value never needed to be
known. A test therefore consists of a modification step, a restoring step, and a comparison.
A key rule is that the instruction under test appears only in the modification step, so that
a systematic fault cannot cancel itself out during restoration.

\subsection{The PATARA Tool Flow}
PATARA is an ISA-independent generator with three stages: instruction selection, sequence
generation, and code generation. Each instruction is described in XML together with its
inverse, so that the generator can be extended to a new (including custom) instruction by
adding a single description rather than modifying the tool. The output is self-checking
assembly that runs through the existing toolchain and simulator and needs no external
reference.

\subsection{Extensions for Pipeline Corner Cases}
Plain REVERSI leaves structural blind spots, which are addressed by three extensions.
Hazard-sequence generation systematically enumerates instruction combinations to exercise
forwarding and stall interactions, raising condition coverage from about 81\% to 97\%.
Operand swapping reaches the second operand's forwarding path, which basic tests never
exercise. Cache-miss generation deliberately forces misses to test miss and stall handling,
closing the remaining gap to full coverage.

\subsection{The Twin Method}
REVERSI cannot verify operations that discard information, such as an absolute value or a
dot product, because they cannot be inverted. For these, the same operation is computed on
the hardware under test and, separately, as software on the already-verified core; the two
results are then compared. Because the core is verified first by REVERSI, it can serve as
the trusted reference for the co-processor, so no external simulator is required. The
twin-based framework is shown in Figure~\ref{fig:twin}.

\begin{figure}[h]
\centering
\figplaceholder{Twin-based verification framework (Fig.~7 of the base paper)}
\caption{Twin-based verification framework of PATARA. Source: Gesper et al. (2026).}
\label{fig:twin}
\end{figure}

\section{System Architecture/Framework}
As shown in Figure~\ref{fig:twin}, a constrained-random generator produces an instruction
sequence from the instruction descriptions. The sequence is executed both on the
co-processor (the hardware twin) and, in parallel, as a software reference on the verified
RISC-V core (the software twin). The two results are compared during a validation stage to
produce the verification result. Constraints on the generator---such as valid vector-lane
chaining and the avoidance of hazards on the vertical vector unit---ensure that only legal
instruction sequences are generated. Because the verified RISC-V core itself provides the
reference, the entire system checks itself from the bottom up, without any external golden
model.

\section{Experimental Setup}
The framework was evaluated on a custom six-stage RV32IM pipeline with the stages IF, ID,
EX, MEM1, MEM2, and WB. The execute stage contains separate single-cycle ALU and
multi-cycle multiply and divide units, together with full forwarding logic; the design is
therefore representative of a real pipelined core in which timing-dependent bugs can arise.
The attached co-processor is the V\textsuperscript{2}PRO vertical vector unit.

The generated self-checking tests were compiled through the standard RISC-V toolchain and
executed on the register-transfer-level simulation of the core. Coverage was measured using
a commercial simulator across six metrics---statement, branch, expression, condition,
finite-state-machine transition, and toggle coverage---and the results were compared against
those obtained by running the official RISC-V compliance suite on the same design. Condition
coverage, the strictest of these metrics, was used as the primary basis for comparison.

\section{Results and Performance Analysis}
The official compliance suite executed on the core without detecting any bug, yet left
condition coverage at 79.12\%. PATARA reached 100\% condition coverage and revealed several
real bugs in the same design, all of them timing or interaction faults in control
logic---for example, forwarding during a multi-cycle divide, a cache miss during a multiply,
and incorrect forwarding from a load into a JALR instruction.

A notable observation is that PATARA's simplest configuration already exceeded the
compliance suite's coverage while using roughly 127 times fewer instructions, and that full
coverage was reached only after adding the hazard-sequence and cache-miss extensions.
Table~\ref{tab:results} shows how condition coverage improves as each capability is added.
On the co-processor, the twin method also reached 100\% coverage; a tenfold increase in the
number of random tests produced no improvement, and the final gap was closed only by biasing
generation toward short vectors---illustrating that test \emph{quality}, not quantity, drives
coverage.

\begin{table}[h]
\centering
\caption{Condition Coverage by Test Configuration}
\label{tab:results}
\begin{tabularx}{\textwidth}{Xrr}
\toprule
\textbf{Configuration} & \textbf{Instructions} & \textbf{Condition coverage}\\
\midrule
Official compliance suite            & 889{,}516   & 79.12\%\\
PATARA: instruction combinations     & 6{,}983     & 80.21\%\\
\quad + operand swapping             & 9{,}722     & 81.31\%\\
\quad + hazard sequences             & 1{,}851{,}852 & 96.70\%\\
\quad + cache-miss tests             & 6{,}732{,}901 & 100.00\%\\
\bottomrule
\end{tabularx}\\[4pt]
\footnotesize Source: Prepared by the author based on the results reported in Gesper et al.
(2026).
\end{table}

\section{Comparative Analysis with Existing Methods}
Compared with the golden-model approach of the Vector Accelerator paper, PATARA is
distinguished by requiring no reference model at all: correctness is established by
self-consistency for invertible instructions and by an on-chip software twin for
non-invertible ones. Relative to the formal methods ($\chi$RVFormal and PFV), PATARA is far
more scalable and directly applicable to a real pipelined core, but it provides a weaker
guarantee---it proves self-consistency rather than conformance to the ISA specification. The
study also shows that high coverage and high bug yield are distinct properties: the
golden-model campaign reported far more errors at lower coverage, whereas PATARA reports full
coverage with fewer, self-reported bugs. This trade-off reinforces the central theme that no
single verification method simultaneously maximises completeness, scalability, and
independence from a reference.

\section{Advantages, Limitations, and Future Scope}
\noindent\textbf{Advantages}
\begin{itemize}[nosep]
  \item Requires no golden reference model.
  \item Generates tests that target the specific pipeline of the design under test.
  \item Achieves 100\% condition coverage and finds real bugs missed by the compliance
        suite.
  \item Is extensible to custom instructions through a single XML description.
  \item Produces self-checking tests that remain usable post-silicon.
\end{itemize}

\noindent\textbf{Limitations}
\begin{itemize}[nosep]
  \item Covers RV32IM only, excluding floating point (whose rounding breaks
        invertibility) and the A and C extensions.
  \item Excludes CSRs, interrupts, and the privileged specification.
  \item The twin method can be defeated by an erroneous software twin and does not catch
        purely timing-related bugs.
  \item 100\% coverage does not imply correctness: REVERSI proves self-consistency, not
        conformance to the specification.
  \item The bug evidence is limited to a single, self-reported design.
\end{itemize}

\noindent\textbf{Future Scope}
\begin{itemize}[nosep]
  \item Extending the approach to floating-point and privileged instructions.
  \item Combining self-testing with a formal, specification-derived reference to close
        the self-consistency gap.
  \item Confirming the discovered bugs on third-party cores.
\end{itemize}

\vspace{0.3cm}
\noindent\textbf{Chapter Conclusion}\\[2pt]
This chapter presented a detailed analysis of the selected base paper, which proposes a
self-testing framework for verifying a custom RISC-V core and co-processor without a golden
reference model. The methodology combines the REVERSI self-checking principle, an
XML-driven generation tool, three extensions for pipeline corner cases, and a twin method
for non-invertible operations. The experimental analysis demonstrated that the compliance
suite passed a design that in fact contained real bugs, whereas the proposed framework
reached full condition coverage and exposed those bugs. Although the approach provides a
weaker guarantee than formal verification, its independence from a reference model and its
support for custom instructions make it directly relevant to the proposed project, which is
discussed in the following chapter.

%---------------------------------------------------------------------
\chapter{Project Relevance and Conclusion}

\section{Summary of Key Findings}
The literature review and the detailed analysis of the base paper provided valuable insight
into how correctness is established for RISC-V processors. The reviewed papers demonstrated
that verification methods differ along two axes---where test stimuli originate and how
correctness is decided---and that no single method is simultaneously complete, scalable, and
reference-free. Formal methods give the strongest guarantees but either bound their proofs
to a fixed depth or apply only to trivial cores; golden-model testing and fuzzing scale to
real designs but check only sampled inputs and depend on a reference for the standard ISA.

The base paper was shown to be the only studied method that verifies a processor without any
golden model. It establishes correctness through the REVERSI self-checking principle for
invertible instructions and through a twin method for non-invertible operations, and it
demonstrates concretely that passing the compliance suite is not the same as being correct.
These findings provide a strong technical foundation for the verification of custom
instructions, which is the central concern of the proposed project.

\section{Relevance to the Proposed Project}
The concepts studied during this seminar have direct relevance to the author's project,
which implements custom AI-acceleration instructions on the Shakti C-Class core. Because
these instructions are invented by the designer, no standard golden reference---such as the
Spike simulator---recognises them, which is exactly the problem the base paper addresses.

An invertible custom instruction such as ADDP1 ($rd = rs1 + rs2 + 1$) maps directly onto the
REVERSI principle: its inverse is a short sequence of already-verified standard
instructions, namely a subtraction followed by a decrement. A non-invertible custom
instruction such as PDOT8, an eight-lane dot product, cannot be inverted because the result
discards its inputs, and therefore uses the twin method. In the author's work, PDOT8 was
verified by computing the same result in hardware and in software on the verified core,
producing zero mismatches over 16{,}384 outputs---a direct reproduction of the base paper's
methodology.

\subsection{Contribution of Seminar Concepts to the Project}
The seminar contributes three things to the project. First, it provides a concrete
verification methodology for custom instructions: REVERSI self-tests for invertible
operations and the twin method for non-invertible ones. Second, it establishes an awareness
that the Shakti C-Class pipeline, which uses forwarding and shared decode and ALU logic that
the project modifies, is exposed to precisely the timing-interaction bugs the base paper
targets. Third, it provides a stronger justification for correctness than the current
practice of relying on the compliance suite.

\section{Proposed Project Direction}
The proposed project is the design and verification of custom instruction-set extensions for
AI acceleration on the Shakti C-Class core. Building on the knowledge gained from the
seminar, each custom instruction will be verified using the methodology studied here:
REVERSI self-tests for invertible operations and the twin method for non-invertible ones.
The compliance-suite-based justification currently used will be replaced by systematic,
self-checking tests that reach the timing-dependent corner cases the compliance suite leaves
untested. Formal, specification-derived checks identified in the supporting papers are noted
as a complementary direction for the standard portions of the core.

\section{Scope for Improvement and Future Work}
Although the base-paper methodology is directly applicable, several opportunities exist for
extension. The verification coverage can be extended to the floating-point and compressed
extensions supported by the Shakti core, which the base paper does not address. The twin
method can be applied to further custom instructions as they are designed. Finally, a
lightweight formal check could be introduced to close the self-consistency gap inherent in
self-testing, thereby combining the scalability of self-testing with the specification
conformance offered by formal verification.

\section{Conclusion}
This seminar provided a comprehensive understanding of verification methodologies for
RISC-V processors. Through the literature review and the detailed analysis of the selected
base paper, the principal approaches---formal model checking, golden-model functional
verification, coverage-guided fuzzing, and polynomial formal verification---were studied,
together with a self-testing framework that verifies a custom core and co-processor without
a golden reference model.

The central lesson of the seminar is that verification is a matter of trade-offs:
completeness, scalability, and independence from a reference cannot all be maximised at
once. For a project built around custom instructions that have no reference, the
self-testing and twin-based approach of the base paper is the most directly applicable, and
it has already been used successfully to verify the project's own custom instruction. By
extending the concepts presented in the base paper, the proposed project aims to contribute
towards the reliable and systematic verification of custom RISC-V instruction-set
extensions for edge AI acceleration.

%=====================================================================
% REFERENCES
%=====================================================================
\begin{thebibliography}{9}
\addcontentsline{toc}{chapter}{References}

\bibitem{patara} S. Gesper, F. Stuckmann, L. W\"obbekind, and G. Pay\'a-Vay\'a,
``A Self-Testing Framework for Verification and Validation of a RISC-V-Based System with a
Co-processor,'' \textit{International Journal of Parallel Programming}, vol.~54, art.~15,
2026.

\bibitem{xrvformal} S. Shen, S. Chen, Y. Liu, L. Zhang, F. Song, and Z. Wu,
``$\chi$RVFormal: Formal Verification of RISC-V Processor Chisel Designs,''
\textit{Journal of Systems Architecture}, vol.~175, art.~103761, 2026.

\bibitem{vector} V. Jim\'enez et al., ``Functional Verification of a RISC-V Vector
Accelerator,'' \textit{IEEE Design \& Test}, vol.~40, no.~3, pp.~36--44, 2023.

\bibitem{instiller} G. Zhang, P. Wang, T. Yue, D. Liu, Y. Guo, and K. Lu,
``INSTILLER: Toward Efficient and Realistic RTL Fuzzing,'' \textit{IEEE Transactions on
Computer-Aided Design of Integrated Circuits and Systems}, vol.~43, no.~7, pp.~2177--2190,
2024.

\bibitem{pfv} L. Weingarten, K. Datta, and R. Drechsler, ``Polynomial Formal Verification
of a RISC-V Processor,'' \textit{IEEE Transactions on Nanotechnology}, vol.~24,
pp.~140--151, 2025.

\bibitem{riscvman} A. Waterman and K. Asanovi\'c, ``The RISC-V Instruction Set Manual,
Volume I: Unprivileged ISA,'' RISC-V International, 2019.

\bibitem{reversi} E. Shazli and M. Tehranipoor, ``REVERSI: Post-Silicon Validation via
Self-Checking Randomized Test Programs,'' foundational REVERSI method extended by the base
paper.

\end{thebibliography}

%=====================================================================
% QUESTIONS AND ANSWERS
%=====================================================================
\chapter*{Questions and Answers}
\addcontentsline{toc}{chapter}{Questions and Answers}

\begin{enumerate}[leftmargin=*,itemsep=10pt]
\item \textbf{Does 100\% coverage mean the processor is correct?}\\
\textbf{Answer:} No. It means every line, branch, and logical condition of the written code
was exercised. REVERSI proves self-consistency---that a subtraction undoes an
addition---rather than conformance to the RISC-V specification, so coverage is necessary but
not sufficient. This is one reason formal-verification papers are also included in the study.

\item \textbf{Why is PATARA chosen as the base paper rather than the formal-verification
paper?}\\
\textbf{Answer:} A base paper should be the one the project builds on. The project reproduced
PATARA's twin method on its own custom instruction. Moreover, formal verification proves
conformance to the specification, but custom instructions are not in the specification, so it
cannot verify the very thing the project needs.

\item \textbf{Could any of the supporting papers verify a custom instruction directly?}\\
\textbf{Answer:} Not off the shelf. Each relies on a reference---a formal model, Spike, or an
instruction-set simulator---that describes only the standard ISA, so the reference would
first have to be extended by hand.

\item \textbf{What did $\chi$RVFormal actually contribute?}\\
\textbf{Answer:} Two things: a compact Chisel reference model, and, more importantly, a
synchronisation mechanism (StateCheck and ActionCheck) that allows a single-cycle reference
model to be compared against a pipelined or out-of-order core.

\item \textbf{How can a single-cycle reference verify an out-of-order core such as BOOM?}\\
\textbf{Answer:} An out-of-order core executes out of order but commits in program order
through its reorder buffer, so the sequence of committed effects matches a single-cycle
machine. The comparison is therefore performed only at instruction commit.
\end{enumerate}

%=====================================================================
% APPENDIX A: PRESENTATION
%=====================================================================
\chapter*{Appendix A: Presentation}
\addcontentsline{toc}{chapter}{Appendix A: Presentation}
% To embed the slides, add \usepackage{pdfpages} to the preamble and use, e.g.:
% \includepdf[pages=-,nup=2x2,frame,scale=0.92]{KEVIN_JOSE_SEMINAR_PRESENTATION.pdf}
\textit{[Insert presentation slides here using the pdfpages package.]}

\end{document}
